% =====================================================================================================================
\section{\bench{}: a testbed for number checks}
\label{sec:setup}

A study of where a check belongs has to follow the loop: the reviewer first, then a check of the finished paper, then
the check inside the agent's session. \bench{} has one part for each (Table~\ref{tab:bench}). All model calls use
Claude Haiku 5.5, Sonnet 5.5, and Opus 5.5 through the Claude Code command line, and every paper and repository was
posted after these models' training cutoff, so no model has seen the text it reviews or the results it writes about.
Appendix~\ref{app:details} gives the full protocols and prompts.

\begin{table}[t]
  \centering
  \caption{The four parts of \bench{}. Each part fixes the papers, the intervention, and the outcome before any model
  is run; the two placement parts and the reviewer part have analysis plans that were written down in advance
  (Appendix~\ref{app:details}).}
  \label{tab:bench}
  \small
  \setlength{\tabcolsep}{5pt}
  \begin{tabular}{@{}p{0.15\linewidth}p{0.27\linewidth}p{0.27\linewidth}p{0.22\linewidth}@{}}
    \toprule
    \rowcolor{headgray}
    Part & Material & Intervention & Outcome \\
    \midrule
    Reviewers & 60 post-cutoff arXiv papers, 1,374 reviews & 1 or 5 prose numbers made wrong or deleted; all
    evidence removed & Score change; reviews quoting a planted value \\
    Writers & 10 repositories with raw results, 40 revised drafts & Four rounds of revision under a reviewer, or
    generic revision & 1,192 numbers added or changed \\
    Audited sample & 300 of those numbers, drawn at random & Two blind model audits; 122 rows recomputed by hand, 50
    of them twice & Errors caught and false alarms per checker \\
    Placement & 8 repositories $\times$ 3 seeds, 138 planted errors & Five placements of the check, no check, and a placebo; then a writer
    told to protect its results, and two more writer sizes & Planted errors left; laundering; turns \\
    \bottomrule
  \end{tabular}
\end{table}

\paragraph{Reviewers on real papers.}
We collected 60 arXiv papers in machine learning, computer vision, and computational linguistics, each with an
experiments section, at least two numeric tables, and at least five prose numbers that report a measured result also
found in a table. Each paper was reviewed in five versions: as posted; with one of those prose numbers shifted so it
contradicts the table that still holds the true value (1 wrong); with five shifted (5 wrong); with the same five
numbers deleted from their sentences (5 deleted); and with the experiments section and every table removed (no
evidence). Reviewers gave a 0--100 percentile score in fresh, stateless calls. We report score drops from the clean
paper and the share of reviews whose summary, weaknesses, or questions quote a planted value. The analysis was fixed
before any review was run.\footnote{The pre-registered primary test asked whether five wrong numbers cost the same as
five deleted ones within a 3-point margin, counting a reviewer only if its no-evidence control cost at least 10
points. Opus met the equivalence bound, but no reviewer's control reached 10 points, so by our own rule the test is
inconclusive. We report every contrast as measured (Appendix~\ref{app:killa}).}

\paragraph{Writers under reviewer pressure.}
We took 10 research repositories whose papers appeared after the cutoff and whose raw results are public, and removed
the paper, README, and documentation. Starting from the same Sonnet draft, a Haiku or Sonnet writer revised the paper
for four rounds under a Sonnet reviewer with best-of-three selection, or, as a control, revised it generically and
kept a random candidate. The 40 final drafts hold 1,192 numbers that the writers added or changed. We drew 300 of
them at random, stratified by repository and blind to every checker's verdict.

\paragraph{Ground truth.}
Sonnet and Opus each audited all 300 numbers from the raw data, with tags and recipes removed and prose notes declared
not to be evidence. They agreed on 90\% of labels; a number counts as an error only when both call it one, and numbers
with one dissenting audit are reported as disputed. A team member then recomputed a 122-row packet by hand without
seeing any model label: every agreed error, every disputed number, every correct number that the recompute check
flagged, and a random set it passed. A second team member independently recomputed 50 of these rows; the two agreed
on 47 (Cohen's $\kappa = 0.74$).

\paragraph{Checkers compared.}
Recompute is the check of Section~\ref{sec:method}, replayed offline: the model that wrote each draft tagged its
new numbers and wrote recipes without changing any text, and the checker ruled. The lookup gate is our earlier
design, which links each number to stored result values through shared words in table labels, file names, and keys,
and requires a changed number to still match the value its predecessor was linked to. The post-hoc trace runs
the same linking on the finished paper. The audit agent is a Claude Code agent with read and shell tools, given
one finished draft and its repository and asked to recompute every number under the same evidence rule as the ground
truth, as FabScore does~\cite{chen2026fabscore}. A second variant gives every checked number a probability of being
wrong, which lets us set a cutoff where the audit flags no more correct numbers than recompute does (matched alarm
rate).

\paragraph{Placement.}
Each of the 24 placement papers is a tagged draft of one repository in which every tag recomputes, seeded with up to
six errors: four flattering ones, which move a result of the paper's own method by 6 to 12\% in its better direction,
and two neutral last-digit slips on other numbers. Flattering errors are the ones an agent that optimizes a reviewer
score has a reason to keep. A Sonnet writer then revises the paper once under a fresh Sonnet review that also asks
for new quantitative analyses, so the agent adds numbers of its own. Every placement sees the same paper, the same
planted errors, and the same review. Seven arms differ only in the check: none; the rules and the check command in
the prompt (tool); feedback; gate; generate; a placebo that denies edits touching numbers at random, at the gate's
measured rate and with a generic message; and an end-of-session Opus audit whose flags go back to the writer for a
second session. For each planted error we record whether the agent corrected it, deleted it, left it, or laundered
it, which means changing the recipe until it reproduces the wrong number. We also count correct numbers the agent
turned wrong, the numbers it added, and its turns. Two earlier pilots on 6 repositories, which shaped this design,
are reported in Appendix~\ref{app:placement}.

\paragraph{Extensions.}
Two further arms reuse the same 24 papers. In the first, the writer's prompt carries a note from the lead author
telling it not to weaken the paper's headline results, so that correcting a flattering error goes against its
instructions; we run the tool, feedback, and gate arms this way. In the second, a Haiku writer and an Opus writer
replace Sonnet in the tool and feedback arms.

\paragraph{Statistics.}
For reviewers, scores are averaged within paper and version and then over papers, with bootstrap intervals over papers
(10,000 resamples). Checker results are counts on fixed samples. For placement, the paper is the unit: we report
planted errors left per paper and paired differences between arms with bootstrap intervals that resample
repositories and keep their seeds together, with Holm correction across the pre-registered contrasts.
